\section{Limitations and research agenda}\label{sec:discussion}

\subsection{Annotation and trust}

An incomplete reasoning graph can mislead precisely because it looks explicit. Missing premises, merged quantifiers, and misclassified evidence can hide behind a clean diagram. Phase II's review experience motivates source anchors and explicit scope, but those features do not establish annotation accuracy automatically. Automatic extraction should therefore produce reviewable drafts, with coverage accounting and targeted human checking. 

The case study combines a selected working graph, an agent-led investigation, written proofs, and exact certificate checks. Its strongest mathematical outcome is the uniform theorem in \cref{thm:main}. The methodological question is broader: whether maintaining this structured history improves later research, review, and learning. That question calls for the comparative and cumulative-use studies in \cref{sec:evaluation}.

\subsection{Scale, selectivity, and failure modes}

A collection can become unusable if every generated suggestion is retained as an undifferentiated node. RRE needs separation between an archival event store and a curated current view. Unreviewed proposals can remain searchable without appearing as supported mathematical facts. Duplicate detection, summary views, and relevance filters should reduce navigation cost while preserving access to the underlying evidence.

Search can also become trapped by the existing graph. If a promising route lies outside the represented literature, repeated local recombination may never find it. The protocol should permit source expansion, alternative representations, and human redirection. Any claim that graph-guided exploration is complete would require assumptions far beyond the present proposal.

The method can create incentives to optimize visible activity rather than mathematics. Provisional significance labels should therefore be justified by resolved obligations and downstream uses, not by the size of the graph or the length of a run. A precise counterexample may be more valuable than many speculative connections. A clear statement that no progress was made is a legitimate outcome.

\subsection{Community concerns and participation}

RRE addresses a specific set of concerns: hidden dependencies, loss of context, duplicated failed work, unclear attribution, and the difficulty of inspecting AI-assisted claims. It does not settle disagreements about commercial AI, environmental costs, labor, or whether particular forms of automation belong in mathematical practice. An inspectable record is one contribution to responsible research, not a reason to dismiss those concerns.

Participation should be voluntary. Researchers should be able to use the infrastructure without AI-generated contributions, restrict automated activity, and decide which records become public. Shared schemas and portable exports can reduce dependence on one provider. Human mathematical judgment remains essential in setting worthwhile objectives, interpreting significance, and deciding what deserves community attention.

\subsection{Next steps}

The immediate technical agenda is to implement a controller that enforces the protocol and supports concurrent contributions, while extending the versioned companion with persistent target management and automatic correction propagation. It should integrate symbolic computation and proof assistants where useful, retain explicit limits when formalization is incomplete, and export both a readable narrative and machine-readable evidence. The proposed schema in \cref{app:records} is a starting point for that implementation rather than a completed interoperability standard.

Open-ended autonomous exploration would extend this controller to selecting overall research objectives from the collection. Research-strategy improvement is a distinct extension: the system could propose changes to retrieval, target selection, conjecture generation, or checking tools, then evaluate those changes before adopting them. Each revision should retain the previous method, evaluation tasks, resource costs, results, and acceptance decision. These records concern the research procedure and belong alongside the mathematical graph; they are not premises proving its theorems. More graph records alone would establish neither better research methods nor greater mathematical value.

For tensor certificates, the next mathematical targets are precise. The nilpotent-pencil argument leaves two possible degrees, $3m+3$ and $3m+6$, for arbitrary nonzero nilpotent types. For three or more summands, the transfer lemma already applies; the next task is to find reductions to base tensors with certificates sharp enough to determine exact degrees. These targets can use the recorded identities and degree ledger directly. For RRE, independent reuse of this package and controlled comparisons provide the next methodological tests.

\section{Conclusion}

The Math Instruments program progresses from identifying reusable devices, through reconstructing their reasoning, to investigating new questions with an evolving mathematical record. RRE proposes a disciplined form of that third phase. Each cycle connects an objective to a target, an investigation to evidence, and an outcome to the next frontier. Supported results, scoped failures, and unresolved obligations all contribute to the collection.

The case study turns an inherited lower bound and a missing boundary argument into an exact theorem for the entire two-term tensor-decomposition model. The reusable transfer rule comes from combining the general balance reduction with explicit degree accounting for contracted modes. Applying it to the degree-9 boundary certificate and then using classical geometry resolves the original infinite-family objective and more. The graph preserves the intermediate pencil theorem, the route to the general result, and the remaining questions. This is the proposed form of cumulative research: mathematical outcomes return to a record that others can inspect, extend, and use to formulate their next investigation.

\section*{Data and software availability}
The local companion package contains the certificate-study charter, 44 exploration records, six graph versions and their deltas, mathematical notes, exact certificates, checking programs, source pins, and validation reports. It also contains the versioned graph viewer and the figure-generation and capture programs. \Cref{app:records} describes replay; \cref{app:complexity-graphs} documents the underlying atlas and its representative sample. The two earlier phases of the Math Instruments program are available through the versioned Zenodo records in \cite{xie1,xie2}. The new package is a local preprint companion; its manifest records the included artifacts and their hashes.

\section*{Acknowledgment}
The author thanks OpenAI's GPT-6.0 Astra for conducting the goal-directed autonomous Recursive Reasoning Exploration (RRE) run under the author's research objective and authorized milestone, including developing the arguments, writing checking programs, generating the follow-up target, and assisting with manuscript drafting. GPT-6.0 Astra was also used to develop the v2 companion visualization of reasoning-graph dynamics. The author established the research direction, designed the RRE methodology, set the case study's objective and authorized milestone, performed an initial audit of the mathematical arguments for correctness, and takes responsibility for the final presentation of the paper.
